%% ===================================================================
\section{\vrev{연구 요약}}
\label{sec:conclusion-summary}
%% ===================================================================

현행 BCMS에서 RTO는 BIA 시점에
고정값으로 설정되며, 이후 운영 환경의 변화와 무관하게 유지된다.
이러한 정적(static) RTO 접근법은 세 가지 구조적 한계를 내포한다.
첫째, 조직의 관측성(observability) 수준이 시간과 상황에 따라 변동함에도
이를 반영하지 못한다.
둘째, 복구 과정에 내재된 불확실성(uncertainty)을 확률 분포에 기반하여
정량화하지 않는다.
셋째, 사이버 공격이나 복합 재난 등 극단적 사건의 파레토(중꼬리, Heavy-Tail) 특성을
구조적으로 포착하지 못한다.

본 연구는 이러한 한계를 극복하기 위하여, 관측성과 불확실성을
개별 시그모이드 바운딩(individual sigmoid bounding) 함수로 통합하는
$\DRTO$ 프레임워크를 제안하였다.
핵심 수식은 \m{식~(\ref{eq:ch7-drto-core})에 나타낸 바와 같이} 정의된다.

\begin{equation}
  \DRTO = R_0 + E_{O,\text{bnd}} + E_{U,\text{bnd}}
  \label{eq:ch7-drto-core}
\end{equation}

\noindent
여기서 관측성 효과 $E_{O,\text{bnd}} = -R_0 \cdot S(z_O)$는
$\DRTO$를 기준선 $R_0$로부터 하향 조정하고,
불확실성 효과 $E_{U,\text{bnd}} = (R_{\max} - R_0) \cdot S(z_U)$는
상향 조정하며, 수정 시그모이드 $S(z) = 2/(1 + e^{-z}) - 1$에 의해
물리적 경계 $\DRTO \in (0,\; R_{\max})$가 점근적으로 보장된다.
$R_0 = 6.0$시간, $R_{\max} = 24.0$시간, $\kappa = 1.2$를 기본 설정으로 채택하였다.

연구 검증은 세 축으로 수행되었다.
첫째, $n = 10{,}000$ 몬테카를로 시뮬레이션을 4개 조건(C0--C3)에 대해
실시하여 \jrev{10개 연구가설}(H1--\jrev{H10})의 모델 내부 정합성을 체계적으로 확인하였다.
둘째, BCM/DR 분야 전문가 5명을 대상으로 구조화된 설문(27문항, 5점 리커트 척도)과
심층 인터뷰(약 30분, 반구조화)를 수행하여 모델의 실무 타당성을 삼각 검증하였다.
셋째, 1차 단순, 1차 다중, 2차 회귀 분석 및 P85.8 기준 분할 회귀를 통해
모델 응답 곡면의 구조적 특성을 규명하고, 꼬리 감쇠 현상을 발견하였다.
그 결과, \textbf{\jrev{10개} 가설 전원이 사전 설정 기준을 충족하여 채택}되었다.
넷째, \chg{\citet{lee2026drto}}과 본 연구의 5개 연구 단계에 대한
\chg{연구모형, 온톨로지 분석, 논문 내용, 기준값} 간
\textbf{303개 교차검증 항목을 전수 확인}하여
\chg{수치, 수식, 용어}의 완전 정합(전원 PASS)을 달성하였으며,
이는 연구 프레임워크의 논리적 일관성과 데이터 신뢰성을 보증한다.


